\section{结束语}
针对动态热点下的多无人机协同计算，建立了区分上传、转发和计算阶段的任务级服务模型，以全部分片完成定义父任务端到端时延，并用未完成父任务驻留量保留跨帧等待信息。在此基础上，将部署与任务路由组织为共享物理环境的大、小时间尺度决策，采用统一物理折扣处理两层反馈和部署空档，给出了PPO与Direct SAC的采样、执行及更新过程。单纯形映射保证路由比例可行，执行路由评论家满足相同执行动作的价值表征一致性；这些性质不直接推出学习策略收敛或队列稳定。后续性能验证需同时覆盖完成时延、未完成任务、积压与能量，并通过独立对比检验两层学习的实际增量价值。
